\section{The truncated hierarchy}\label{app:truncated}

This appendix gives the proofs and derivations behind section~\ref{sec:truncated}, whose notation it uses: $A$ is the $(M+1)\times(M+1)$ ladder matrix of the closure in use, $\Pw$ its weight, $u_j$ its eigenvalues in increasing order, $\boldsymbol r_j$ and $\boldsymbol w_j=\Pw\boldsymbol r_j$ the right and left eigenvectors, $n_j=\boldsymbol r_j^{\mathsf T}\Pw\boldsymbol r_j$, $R_{mj}=r_m(u_j)/\sqrt{n_j}$ and $\lambda_j=1/n_j$. We write $\boldsymbol e_m$ for the $m$th unit vector, which is the coefficient vector of the Hermite basis function $e_m$, $h_m(u)=H_m(u)/\sqrt{2^mm!}$ for the polynomials orthonormal with weight $F_M$, and $\psi_m=h_mF_M^{1/2}$ for the orthonormal Hermite functions. The nodes and weights of the $(M+1)$-point Gauss--Hermite rule, normalised so that the weights add up to one, are $x_j$ and $\lambda^\infty_j$. Each identity below was checked symbolically in $\Lambda$ for small $M$, or numerically, as stated.

\subsection{Weights, closures and eigenvectors}\label{appC:weights}

Let $J$ be the ladder matrix of the zero closure at $\Lambda=\infty$, the symmetric Jacobi matrix of the $h_m$, and $J_c$ that of the copy closure, which for $M\ge2$ differs from $J$ only in the entry $(J_c)_{M,M-1}=\sqrt{M/2}+\sqrt{(M+1)/2}$. A finite $\Lambda$ changes only the entry $(1,0)$, where $A_{10}=(1-1/\Lambda)\,J_{10}$, and column $0$ of $J$ and of $J_c$ has the single non-zero entry $J_{10}=1/\sqrt2$. Hence, with $P_\Lambda=I-\boldsymbol e_0\boldsymbol e_0^{\mathsf T}/\Lambda=\mathrm{diag}(1-1/\Lambda,1,\dots,1)$,
\begin{equation}\label{appC:factor}
  A=J\,P_\Lambda\quad\text{(zero closure)},\qquad A=J_c\,P_\Lambda\quad\text{(copy closure, $M\ge2$)}.
\end{equation}
This is the Hermite form of the factorisation (\ref{eq:LUP}), $L=UP$: $J$ is the matrix of multiplication by $u$, and $P_\Lambda$ that of $P$. For the zero closure $\Pw=P_\Lambda$, so that $\Pw A=P_\Lambda JP_\Lambda$ is symmetric. For the copy closure, let $D_c=\mathrm{diag}(1,\dots,1,p_M)$ with $p_M$ from (\ref{tr:pM}). Then $p_M(J_c)_{M,M-1}=\sqrt{M/2}=J_{M,M-1}$, so that $D_cJ_c=J$, and with $\Pw=D_cP_\Lambda$
\begin{equation*}
  \Pw A=P_\Lambda\,(D_cJ_c)\,P_\Lambda=P_\Lambda JP_\Lambda ,
\end{equation*}
which is the same matrix as for the zero closure. A diagonal matrix $\mathrm{diag}(\pi_m)$ symmetrises $A$ if and only if $\pi_mA_{m,m+1}=\pi_{m+1}A_{m+1,m}$ for every $m$, which fixes each ratio $\pi_{m+1}/\pi_m$, so the weight is unique up to a factor. We checked (\ref{appC:factor}) and the uniqueness symbolically in $\Lambda$ for $M=2$ to $7$ and both closures, and that GANDALF's streaming matrix equals $A$ exactly for $M=2$ to $128$, six values of $\Lambda$ and both closures.

The weight $\Pw$ is positive definite if and only if $1-1/\Lambda>0$, that is, $\Lambda\notin[0,1]$, since $p_M>0$. For $0<\Lambda<1$ the first weight is negative, $\Pw$ is indefinite, and the eigenvalues of $A$ need not be real; the continuum then has the unstable pair of section~\ref{sec:slices}.

The Case--Van Kampen mode $F_u$ satisfies $LF_u=uF_u$, and the ladder matrix without truncation is the matrix of $L$ in the Hermite basis. The Hermite coefficients $r_m(u)$ of $F_u$ therefore satisfy $\sum_nA_{mn}r_n=u\,r_m$ row by row, with $r_0=M_0[F_u]=1$ from the normalisation of $F_u$ in section~\ref{sec:slices} and $r_1=\sqrt2\,M_1[F_u]=\sqrt2\,u$ from row 0:
\begin{equation}\label{appC:rec}
  \begin{gathered}
  r_0=1,\qquad r_1=\sqrt2\,u,\qquad r_2=u\,r_1-\frac{1-1/\Lambda}{\sqrt2}\,r_0,\\
  \sqrt{\frac{m+1}{2}}\;r_{m+1}=u\,r_m-\sqrt{\frac m2}\;r_{m-1}\quad(m\ge2).
  \end{gathered}
\end{equation}
Thus $r_m$ is a polynomial of degree $m$ with the parity $r_m(-u)=(-1)^mr_m(u)$, and $r_m=h_m$ at $\Lambda=\infty$. Explicitly,
\begin{equation*}
  r_m(u)=\frac{(1-1/\Lambda)\,H_m(u)+Y_m(u)/\Lambda}{\sqrt{2^mm!}},\qquad
  Y_m(u)=\int\frac{vH_m(v)-uH_m(u)}{v-u}\,F_M(v)\,\dd v ,
\end{equation*}
because the Dawson-function parts of the principal-value integral in $F_u$ cancel against $\eps_R(u)$; for example, $r_2=\sqrt2\,(2\Lambda u^2-\Lambda+1)/(2\Lambda)$. We checked the recurrence against the Hermite coefficients of $F_u$ computed by principal-value quadrature for $m\le8$, to $9\times10^{-16}$, and the closed form exactly for $m\le10$. Rows $0$ to $M-1$ of the truncated matrix coincide with those of the infinite matrix for both closures, so $(r_0(u),\dots,r_M(u))^{\mathsf T}$ satisfies them for every $u$, and row $M$ is the condition on the eigenvalues,
\begin{equation}\label{appC:eigcond}
  r_{M+1}(u_j)=0\quad\text{(zero closure)},\qquad r_{M+1}(u_j)=r_{M-1}(u_j)\quad\text{(copy closure)}.
\end{equation}
The second condition follows from row $M$ of the copy closure, $u\,r_M=[\sqrt{M/2}+\sqrt{(M+1)/2}\,]\,r_{M-1}$, together with the recurrence at $m=M$. We checked (\ref{appC:eigcond}) and the eigenvectors at $M=16$, 33 and 64, $\Lambda=\pm\sqrt5$ and $4$, and both closures, to $2\times10^{-14}$. The relations $\boldsymbol w_j^{\mathsf T}A=u_j\boldsymbol w_j^{\mathsf T}$, $\boldsymbol w_j^{\mathsf T}\boldsymbol r_l=n_j\delta_{jl}$ and $\sum_j\boldsymbol r_j\boldsymbol w_j^{\mathsf T}/n_j=I$ follow from the symmetry of $\Pw A$, as in the proof of theorem~\ref{thm:trunc}.

\begin{remark}[The copy closure]\label{appC:copy}
GANDALF's closure $g_{M+1}=g_{M-1}$ differs from the zero closure in the following ways. (a) Its eigenvalues are the zeros of $r_{M+1}-r_{M-1}$, so it does not give the Gauss rule of proposition~\ref{tr:gaussprop}, and its largest eigenvalue is larger (table~\ref{tr:umaxtab}). (b) Its weight on $g_M^2$ is $p_M\approx\frac12$, so its conserved free energy is $\sum_{m<M}p_mg_m^2+p_Mg_M^2$, which is not CMM's (B31) on the retained moments, and the diagnostic \texttt{hermite\_free\_energy}, which uses the weight $1$ on $g_M^2$, is not conserved. (c) The matrix $\Pw A$ is the same as for the zero closure, so $Q^{(M)}_u$ is CMM's (B32) for both closures, and the Christoffel--Darboux form of the flux in section~\ref{sec:decorrelation} holds with the eigenvectors normalised by the copy weight. (d) At odd $M$, $A^{-1}=P_\Lambda^{-1}J_c^{-1}$ and $J_c^{-1}=J^{-1}D_c$, so that the conserved helicity is $g^{\mathsf T}\Pw A^{-1}g=g^{\mathsf T}D_cJ^{-1}D_cg$; its matrix differs from the matrix $B_M$ of (\ref{eq:Hoddeven}) by 33\% at $M=9$ and 23\% at $M=33$, relative to the largest element. (e) At even $M$ the left null vector is $\boldsymbol h(0)$ with its last component multiplied by $p_M$, so the energy of the slice at $u=0$ differs from the imaginary part of the helicity of the truncated state through the weight $p_M$ on $g_M$. (f) Both closures are conservative, so both reflect phase-mixing packets (section~\ref{tr:reflsec}).
\end{remark}

\subsection{Gauss quadrature}\label{appC:gauss}

The proof of proposition~\ref{tr:gaussprop} uses the orthonormal polynomials $\sqrt{p_m}\,r_m$ of the measure $\varrho\,\dd u$. Their recurrence follows from (\ref{appC:rec}) and has the Jacobi matrix $\Pw^{1/2}A\,\Pw^{-1/2}$, whose off-diagonal entries are $[(1-1/\Lambda)/2]^{1/2}$ and $[(m+1)/2]^{1/2}$ for $m\ge1$. Truncation with the zero closure keeps its first $M+1$ rows and columns. The weight of the Gauss rule at the node $u_j$ is the value of the Christoffel function there,
\begin{equation*}
  \lambda_j=\Big[\sum_{m=0}^Mp_m\,r_m(u_j)^2\Big]^{-1}=\frac{1}{n_j}=R_{0j}^2 ,
\end{equation*}
which equals $\mu_0O_{0j}^2$ in the form of \citet{GolubWelsch1969}, with $O$ the orthonormal eigenvectors of the truncated Jacobi matrix and $\mu_0=\int\varrho\,\dd u=1/p_0$. The moments of the measure follow from the ladder matrix, $\int\varrho\,u^k\,\dd u=(A^k)_{00}/p_0$. We checked these statements numerically: the eigenvalues of $A$ equal the zeros of $r_{M+1}$ to $3\times10^{-14}$ for $M=10$, 20 and 40, the eigenvectors equal $(r_m(u_j))$ to $9\times10^{-14}$, the weights agree to $5\times10^{-15}$, the moment formula holds for $k\le12$ to $8\times10^{-16}$, and the rule (\ref{eq:gauss}) is exact for all degrees up to $2M+1$ in 50-digit arithmetic at $M=6$ and $10$, to $6\times10^{-44}$. Proposition~\ref{tr:gaussprop}(iii), $\boldsymbol w_j^{\mathsf T}g=\Tr[G](u_j)$ for $G$ in the span of the retained moments, holds to $6\times10^{-16}$ at $M=12$ for random $g$. By the parity of $r_m$, $r_{M+1}(0)=0$ for even $M$, which is the node at $u=0$ of section~\ref{tr:helsec}.

\subsection{Characteristic polynomial and the shift of the phase velocities}\label{appC:shift}

Write $A^{(k)}$ for $A$ with its first $k$ rows and columns removed. Expanding $\det(z-A)$ along its first row,
\begin{equation*}
  \det(z-A)=z\det\big(z-A^{(1)}\big)-A_{01}A_{10}\det\big(z-A^{(2)}\big).
\end{equation*}
Neither $A^{(1)}$ nor $A^{(2)}$ contains the entry $A_{10}$, so both equal the corresponding blocks of the matrix at $\Lambda=\infty$ with the same closure, while $A_{01}A_{10}=(1-1/\Lambda)/2$ and $J_{01}J_{10}=\frac12$. Subtracting the same expansion at $\Lambda=\infty$ gives (\ref{tr:det}), for both closures, since the closure enters only the blocks. We checked (\ref{tr:det}) symbolically in $\Lambda$ and $z$ for $M=3$ to $7$ and in 50-digit arithmetic at $M=30$ and $31$. The traces give the sum rules of section~\ref{tr:velsec}: $\mathrm{tr}\,A=0$, and
\begin{equation*}
  \mathrm{tr}\,A^2=2\sum_{m=0}^{M-1}A_{m,m+1}A_{m+1,m}=\Big(1-\frac1\Lambda\Big)+\sum_{m=1}^{M-1}(m+1)=\frac{M(M+1)}{2}-\frac1\Lambda ,
\end{equation*}
to which the copy closure adds $2\sqrt{M/2}\,\sqrt{(M+1)/2}=\sqrt{M(M+1)}$.

For the zero closure $A=J+\delta A$ holds exactly, with $\delta A=-(\Lambda\sqrt2)^{-1}\boldsymbol e_1\boldsymbol e_0^{\mathsf T}$. The normalised eigenvectors of the symmetric matrix $J$ have components $(\lambda^\infty_j)^{1/2}h_m(x_j)$, so to first order in $1/\Lambda$ perturbation theory gives
\begin{equation}\label{appC:first}
  u_j=x_j-\frac{\lambda^\infty_j\,x_j}{\Lambda}+O(\Lambda^{-2}),
\end{equation}
since $h_0=1$ and $h_1(x)=\sqrt2\,x$. At $M=128$ the remainder is $9.8\times10^{-6}$, $9.7\times10^{-8}$ and $9.7\times10^{-10}$ at $\Lambda=50$, 500 and 5000.

The spectral-shift law (\ref{tr:shift}) can be motivated as follows. By the matrix determinant lemma,
\begin{equation*}
  \frac{\det(z-A)}{\det(z-J)}=1+\frac{\boldsymbol e_0^{\mathsf T}(z-J)^{-1}\boldsymbol e_1}{\Lambda\sqrt2}=\prod_j\frac{z-u_j}{z-x_j}.
\end{equation*}
For $\mathrm{Im}\,z>0$ and $M\to\infty$, the element $\boldsymbol e_0^{\mathsf T}(z-J)^{-1}\boldsymbol e_1$ tends to $\int\psi_0\psi_1\,(z-x)^{-1}\,\dd x=-\sqrt2\,[1+zZ(z)]$, so the ratio tends to $\eps(z)$. We checked this at $z=0.3+0.5\ii$, $-1.2+0.3\ii$ and $2+\ii$, where the differences at $M=512$ are at most $7\times10^{-13}$, $2\times10^{-9}$ and $10^{-13}$ for $\Lambda=\pm\sqrt5$. Suppose now that $u_j-x_j=\theta(x_j)\,\Delta_j$, with $\theta$ varying slowly from node to node. For $\mathrm{Im}\,z$ large compared with the spacing but small compared with the scale of $\theta$, the logarithm of the product is approximately $-\sum_j\theta(x_j)\Delta_j/(z-x_j)\simeq-\int\theta(x')\,(z-x')^{-1}\,\dd x'$, whose boundary value on the real axis is $-\PVint\theta(x')\,(x-x')^{-1}\dd x'+\ii\pi\theta(x)$. The function $\log\eps$ is analytic in the upper half plane and continuous up to the real axis when $\eps$ has no zeros in the closed upper half plane, which is the case for $\Lambda\notin[0,1]$, and it tends to zero at infinity. Matching the two gives $\theta=\pi^{-1}\arg\eps(x+\ii0)$, which is (\ref{tr:shift}), together with $\log|\eps|=-\PVint\theta(x')\,(x-x')^{-1}\dd x'$. This argument assumes the slow variation of $\theta$ and the interchange of the limits, and it does not give the rate at which (\ref{tr:shift}) is approached; the rate of about $M^{-1/2}$ quoted in section~\ref{tr:velsec} is measured. For large $\Lambda$, $\arg\eps\simeq\mathrm{Im}\,\eps=-\pi xF_M(x)/\Lambda$, and with $\lambda^\infty_j\simeq F_M(x_j)\Delta_j$ the law reduces to (\ref{appC:first}).

\subsection{The largest phase velocity}\label{appC:umax}

The normalised extreme eigenvector of $J$ has the components $(\lambda^\infty_M)^{1/2}h_m(x_M)$, so its components on the moments $m=0$ and $1$, on which $\Lambda$ acts, are $(\lambda^\infty_M)^{1/2}$ and $(2\lambda^\infty_M)^{1/2}x_M$. By (\ref{appC:first}) the shift of the largest eigenvalue is $-\lambda^\infty_Mx_M/\Lambda$, with a relative correction of about $1/(2\Lambda x_M^2)$ from second order. In 80-digit arithmetic at $M=12$ and $\Lambda=\sqrt5$ the shift is $-4.921\times10^{-8}$, against $-4.994\times10^{-8}$ from (\ref{appC:first}) (relative difference $0.0145$, close to $1/(2\Lambda x_M^2)=0.013$), and at $M=20$ it is $-5.17\times10^{-14}$. The weight $\lambda^\infty_M$ is of order $e^{-x_M^2}$ up to an algebraic factor: $\ln\lambda^\infty_M+x_M^2$ lies between $-0.60$ and $-0.84$ for $M$ from 12 to 48, while $\ln\lambda^\infty_M+2M+3$ grows from $9.6$ to $15.7$ over the same range, as expected from $x_M^2=2M+3-2^{2/3}|\alpha_1|(2M+3)^{1/3}+\dots$. At $M=64$ the shift is about $10^{-49}$. It is therefore negligible, and $u_{\max}$ is the largest zero of $\psi_{M+1}$ for the zero closure and the largest root of $\psi_{M+1}=\psi_{M-1}$ for the copy closure. We confirmed both by root-finding without an eigenvalue solver.

For the zero closure the first two terms of (\ref{eq:umax}) are the classical asymptotics of the largest zero of a Hermite polynomial. We give the derivation because the copy closure follows the same steps. The function $\psi_{M+1}$ solves $\psi''+(\kappa^2-x^2)\psi=0$ with $\kappa=\kappa_{M+1}=(2M+3)^{1/2}$. With $x=\kappa+(2\kappa)^{-1/3}\xi$ and $\sigma=(2\kappa)^{-4/3}$ the equation becomes, exactly,
\begin{equation*}
  \psi_{\xi\xi}=(\xi+\sigma\xi^2)\,\psi .
\end{equation*}
The uniform approximation $\psi=(\chi')^{-1/2}\mathrm{Ai}(\chi)$, with $\chi(\xi)$ defined by $\chi\,(\chi')^2=\xi+\sigma\xi^2$, gives $\chi=\xi+\sigma\xi^2/5+O(\sigma^2)$, and the term $(\chi')^{1/2}[(\chi')^{-1/2}]''$ that it neglects is $O(\sigma^2)$. The largest zero is at $\chi=\alpha_1$, the first zero of the Airy function, so $\xi=\alpha_1-\sigma\alpha_1^2/5+O(\sigma^2)$, which gives the first line of (\ref{eq:umax}).

For the copy closure at $\Lambda=\infty$ the right eigenvector is $(h_m(u))$, and row $M$ requires $\psi_{M+1}(u)=\psi_{M-1}(u)$. The Hermite functions satisfy $x\psi_M=\sqrt{M/2}\,\psi_{M-1}+\sqrt{(M+1)/2}\,\psi_{M+1}$ and $\psi_M'=\sqrt{M/2}\,\psi_{M-1}-\sqrt{(M+1)/2}\,\psi_{M+1}$, so this requirement becomes the Robin condition
\begin{equation*}
  \psi_M'+\frac{x}{(\sqrt M+\sqrt{M+1})^2}\,\psi_M=0
\end{equation*}
for the function $\psi_M$, which solves the same equation with $\kappa=\kappa_M=(2M+1)^{1/2}$. The coefficient of $\psi_M$ is $1/(2\kappa)$ to leading order, so in the variable $\xi$ the condition reads $\psi_\xi+\sigma\psi=0$. With the same uniform approximation, $\psi_\xi\propto\mathrm{Ai}'(\chi)-[\chi''/(2\chi'^2)]\,\mathrm{Ai}(\chi)$ with $\chi''=2\sigma/5$, so the condition is $\mathrm{Ai}'(\chi)+(4\sigma/5)\,\mathrm{Ai}(\chi)=0$. Near the first zero $\alpha_1'$ of $\mathrm{Ai}'$, where $\mathrm{Ai}''(\alpha_1')=\alpha_1'\mathrm{Ai}(\alpha_1')$, this gives $\chi=\alpha_1'-4\sigma/(5\alpha_1')$, and then $\xi=\chi-\sigma\chi^2/5$ gives the second line of (\ref{eq:umax}).

We computed $u_{\max}$ up to $M=16384$. At $M=16384$ the residual of the two-term formulas times $\kappa^{5/3}$ is $-0.34455$ for the zero closure and $+0.18181$ for the copy closure, approaching the derived $-0.344383$ and $+0.181950$. An unconstrained least-squares fit for $M\ge64$ returns $-1.855757\,\kappa^{-1/3}-0.344271\,\kappa^{-5/3}$ and $-0.808617\,\kappa^{-1/3}+0.182014\,\kappa^{-5/3}$. The residual of the three-term formulas times $\kappa^3$ settles at $-0.1688$ and $-0.1439$, so the remainder is $O(\kappa^{-3})$; these two coefficients are fitted, not derived.

\subsection{The phase-space helicity at odd and even \texorpdfstring{$M$}{M}}\label{appC:hel}

For $m,n\le M$ the matrix of the helicity kernel on the retained moments is
\begin{equation*}
  (B_M)_{mn}=\PVint\frac{\psi_m\psi_n}{u}\,\dd u=\int\frac{h_m(u)h_n(u)-h_m(0)h_n(0)}{u}\,F_M(u)\,\dd u ,
\end{equation*}
because $\PVint F_M\,u^{-1}\,\dd u=0$. The integrand $f_{mn}(u)=[h_mh_n-h_m(0)h_n(0)]/u$ is a polynomial of degree $m+n-1\le2M-1$, so the $(M+1)$-point Gauss--Hermite rule integrates it exactly: $(B_M)_{mn}=\sum_j\lambda^\infty_jf_{mn}(x_j)$. The nodes come in pairs $\pm x_j$ with equal weights, so $\sum_{x_j\neq0}\lambda^\infty_j/x_j=0$. The symmetric matrix $J$ is $O_\infty\,\mathrm{diag}(x_j)\,O_\infty^{\mathsf T}$, with the orthogonal matrix $(O_\infty)_{mj}=(\lambda^\infty_j)^{1/2}h_m(x_j)$.

For odd $M$ no node lies at $u=0$, and
\begin{equation*}
  (B_M)_{mn}=\sum_j\frac{\lambda^\infty_j\,h_m(x_j)\,h_n(x_j)}{x_j}=(J^{-1})_{mn}.
\end{equation*}
By (\ref{appC:factor}), $A^{-1}=P_\Lambda^{-1}J^{-1}$ for the zero closure, so that $\Pw A^{-1}=J^{-1}$ for every $\Lambda\neq1$. The matrix $A$ is invertible at odd $M$ because
\begin{equation*}
  \det A=(-1)^{(M+1)/2}\prod_{i=0}^{(M-1)/2}A_{2i,2i+1}\,A_{2i+1,2i},
\end{equation*}
which vanishes only at $\Lambda=1$; for example, $\det A=105(\Lambda-1)/(16\Lambda)$ at $M=7$ for the zero closure.

For even $M$ the node $x_{M/2}=0$ contributes $\lambda^\infty_{M/2}f_{mn}(0)=\lambda^\infty_{M/2}[h_m'(0)h_n(0)+h_m(0)h_n'(0)]$, and the other nodes contribute $\sum_{j\neq M/2}\lambda^\infty_jh_m(x_j)h_n(x_j)/x_j=(J^+)_{mn}$. This gives the second relation of (\ref{eq:Hoddeven}). The vector $\boldsymbol h'(0)$ has non-zero components only at odd $m$ and $\boldsymbol h(0)$ only at even $m$, so the two are orthogonal, and the rank-two term has non-zero elements between $\boldsymbol h(0)$, the null vector of $J$, and the other eigenvectors of $J$. Hence $B_M$ is not a function of $J$, and at $\Lambda=\infty$ it is not of the form $\Pw\varphi(A)$. The conserved form $\Pw A^D$, with the Drazin inverse $A^D=R\,\mathrm{diag}(\phi_j)\,R^{-1}$, $\phi_j=1/u_j$ for $u_j\neq0$ and $\phi_{M/2}=0$, reduces to $J^+$ at $\Lambda=\infty$, so there it differs from $B_M$ by the node term, $g^{\mathsf T}(\Pw A^D-B_M)g=-2\lambda^\infty_{M/2}\,[\boldsymbol h(0)\cdot g]\,[\boldsymbol h'(0)\cdot g]$. For large even $M$, $\lambda^\infty_{M/2}\simeq F_M(0)\,\pi/(2M+3)^{1/2}=[\pi/(2M+3)]^{1/2}$.

We checked (\ref{eq:Hoddeven}) exactly in SymPy for $M=2$ to $7$, with $\lambda^\infty_{M/2}=2/3$, $8/15$ and $16/35$ at $M=2$, 4 and 6; by direct principal-value quadrature at $M=6$, 7, 10 and 11, to $7\times10^{-16}$, both relations, the odd-$M$ one for $\Lambda=\pm\sqrt5$ and $3$; and the odd-$M$ relation numerically at $M=33$, 65 and 129 for $\Lambda=\pm\sqrt5$, to $7\times10^{-15}$. At $\Lambda=\infty$ the error of the even-$M$ form for $G=F_M(u-\frac12)$ equals the node term to better than $10^{-9}$ relative from $M=16$ on. At finite $\Lambda$ we measured the error only. A fit of its product with $M^{1/2}$ by $c+c'M^{-1/2}+c''M^{-1}$ over $M=64$ to 2048 gives $c=-2.861$ at $\Lambda=\sqrt5$ and $c=-1.008$ at $\Lambda=-\sqrt5$, with residuals below $4\times10^{-5}$.

For the zero closure and even $M$, the vector $\boldsymbol h(0)$ is a left null vector of $A$ for every $\Lambda$. Column $0$ of $\boldsymbol h(0)^{\mathsf T}A$ is $h_1(0)A_{10}=0$, whatever $\Lambda$. The other columns $n<M$ vanish by the recurrence $\sqrt{(n+1)/2}\,h_{n+1}(0)+\sqrt{n/2}\,h_{n-1}(0)=0$, which at $n=1$ involves only $A_{01}$ and $A_{21}$. The last column is $h_{M-1}(0)\sqrt{M/2}$, which vanishes because $h_{M-1}$ is odd. For odd $M$ this entry is not zero (for the vector $(\psi_m(0))$ it is $0.833$ at $M=9$), so that vector is not a null vector; $u=0$ is not an eigenvalue at odd $M$ because $\det A\neq0$, as shown above. The polynomials $r_m$ and $h_m$ obey the same recurrence for $m\ge2$, and $r_2(0)=p_0h_2(0)$, so $r_m(0)=p_0h_m(0)$ for $m\ge1$ and $\boldsymbol w_{M/2}=\Pw\boldsymbol r_{M/2}=p_0\boldsymbol h(0)$. Hence $\boldsymbol w_{M/2}^{\mathsf T}g=p_0\,G_M(0)/F_M(0)=\Tr[G_M](0)$, the continuum value of the slice transform at $u=0$ (section~\ref{sec:slices}), and $a_{M/2}^2=\lambda_{M/2}\,p_0^2\,G_M(0)^2/F_M(0)^2$, whose volume average is proportional to the imaginary part of (\ref{eq:Hcomplex}) for the truncated state. We checked on GANDALF's right-hand side, at $M=10$, 16 and 64, for both signs of $\Lambda$ and both closures, that this slice obeys $\partial_ta_{M/2}+\{\Phi,a_{M/2}\}=0$ to $2\times10^{-15}$ and that $\sum_{\boldsymbol k}|a_{M/2}|^2$ is conserved to relative rates of $2\times10^{-16}$.

\subsection{Closed systems: the time-averaged spectrum}\label{appC:closed}

Consider a single Fourier mode in a uniform field, without flow, collisions or forcing. Each slice then rotates at its own frequency, $a_j(t)=a_j(0)\,e^{-\ii sk_zu_jt}$. With $C_{jl}=a_ja_l^*$, the Hermite spectrum is exactly
\begin{equation*}
  W(m)=p_m|g_m|^2=p_m\sum_{j,l}R_{mj}R_{ml}\,C_{jl},\qquad C_{jl}(t)=C_{jl}(0)\,e^{-\ii sk_z(u_j-u_l)t},
\end{equation*}
which is the truncated form of (\ref{eq:dictionary}). The $u_j$ are distinct, so the average over a long time removes the terms with $j\neq l$ and leaves the slice-diagonal spectrum
\begin{equation}\label{appC:closedspec}
  \overline{W(m)}=p_m\sum_jR_{mj}^2\,E_j,\qquad E_j=|a_j|^2 .
\end{equation}
At $\Lambda=\infty$, $R_{mj}^2=\lambda^\infty_jh_m(x_j)^2\simeq\Delta_j\,\psi_m(x_j)^2$, and for $|x_j|\ll(2m)^{1/2}$ the oscillating factor $\psi_m^2$ averages to $1/[\pi(2m)^{1/2}]$ over neighbouring nodes. The level of (\ref{appC:closedspec}) at large $m$ is therefore $[\sum_jE_j\Delta_j]/[\pi(2m)^{1/2}]$, the discrete counterpart of the level $[\int\mathcal E\,\dd u]/[\pi(2m)^{1/2}]$ of $\delta$-correlated slices in section~\ref{sec:decorrelation}. For a state that starts in $g_0$, $E_j=p_0^2\lambda_j|g_0|^2$, and $\sum_jE_j=p_0|g_0|^2$ is the total free energy $W_{\mathrm{tot}}$. With $\Delta_j\simeq\pi/(2M)^{1/2}$ for the slices with $|u_j|\ll(2M)^{1/2}$, which carry nearly all of it, the time-averaged state is the recurrence equilibrium
\begin{equation*}
  \overline{W(m)}\simeq\frac{W_{\mathrm{tot}}}{2(Mm)^{1/2}} .
\end{equation*}
We tested this at $M=128$ with $sk_z=2\pi$, starting from $g_0$. The time average over $0\le t\le3000$ matches (\ref{appC:closedspec}) to $9\times10^{-5}$, $1.4\times10^{-4}$ and $2.1\times10^{-4}$ of its maximum for $\Lambda=\infty$, $\sqrt5$ and $-\sqrt5$. For $8\le m\le48$, $\overline{W(m)}\,\pi(2m)^{1/2}/\sum_jE_j\Delta_j$ is $0.985$ to $0.997$, $0.973$ to $1.001$ and $0.991$ to $0.999$, and $\overline{W(m)}\,2(Mm)^{1/2}/W_{\mathrm{tot}}$ is $0.98$ to $0.99$, $0.94$ to $0.96$ and $1.00$, for the three values of $\Lambda$. The closed recurrence state and the forced state with nothing returning from $m=M$, whose level is $\epsilon_W/[s|k_z|(2m)^{1/2}]$ (section~\ref{sec:decorrelation}), both have the slope $-1/2$. They differ in the prefactor, and in the flux, which vanishes on average in the closed state.

\subsection{Reflection at the closure}\label{appC:refl}

In the ray derivation, for $\Lambda=\infty$ and $G(0)=F_M$, free streaming at a single wavenumber gives $g_m(t)=(-\ii\varkappa)^me^{-\varkappa^2/4}/(2^mm!)^{1/2}$ with $\varkappa=kst$, so that $|g_m|^2=e^{-\bar m}\,\bar m^m/m!$ is a Poisson distribution in $m$ with mean $\bar m=\varkappa^2/2$. We checked this on GANDALF's linear operator at $M=256$, before the packet reaches $m=M$, to $6\times10^{-15}$. The packet reaches $m=M$ at $t=t_w=(2M)^{1/2}/(|k|s)$. On the way it is damped by the hypercollisions at the local rate $\nu(\bar m/M)^n=\nu(t/t_w)^{2n}$, which reduces its amplitude by the factor $\exp[-\nu t_w/(2n+1)]$. Both closures are conservative and reflect the packet without loss, and the return costs the same factor, which gives (\ref{eq:reflection}) to leading order in $M$.

In the mode derivation the packet is a superposition of the slices near $u=0$. For $|u_j|\ll(2M)^{1/2}$ the eigenvector spreads its weight over $0\le m\le M$ as $R_{mj}^2\simeq\Delta_j/[\pi(2m)^{1/2}]$ with $\Delta_j\simeq\pi/(2M)^{1/2}$, that is, $R_{mj}^2\simeq1/[2(Mm)^{1/2}]$, which gives the damping rate (\ref{tr:gamma}). At $\Lambda=\infty$ the slices near $u=0$ are spaced by $\Delta_0\simeq\pi/(2M+3)^{1/2}$ and re-phase after $T_{\mathrm{rec}}=2\pi/(|k|s\Delta_0)=2t_w[(2M+3)/2M]^{1/2}$, which is $1.0058\times2t_w$ at $M=128$; then $e^{-\mathsf C_{jj}T_{\mathrm{rec}}}$ reproduces (\ref{eq:reflection}). For the six slices nearest $u=0$, $(2n+1)\mathsf C_{jj}/\nu$ is $1.089$, $1.022$, $1.005$ and $1.003$ at $M=64$, 256, 1024 and 2048 for $n=6$ and $\Lambda=\infty$, an excess that falls as $1/M$, and $1.049$, $0.981$, $0.978$ and $0.982$ for $\Lambda=\sqrt5$. The first-order sum over the slices, with the exact $u_j$ and $\mathsf C_{jj}$, reproduces the exact exponent to $1.4\%$ for $M\ge128$ and $X<1.5$, which confirms this mechanism.

At finite $\Lambda$ the spacing at $u=0$ changes (section~\ref{tr:velsec}), and $2\pi/(|k|s\Delta_0)$ would be $1.087\times2t_w$ for $\Lambda=\sqrt5$ and $0.972\times2t_w$ for $\Lambda=-\sqrt5$ at $M=128$. The revival does not occur at these times. Its peak lies at $1.0325$, $1.0049$ and $0.9910$ times $2t_w$ for $\Lambda=\sqrt5$, $\infty$ and $-\sqrt5$ at $M=128$, and at $1.0149$, $1.0012$ and $0.9944$ times $2t_w$ at $M=512$, against $[(2M+3)/2M]^{1/2}=1.0058$ and $1.0015$. The revival time therefore tends to $2t_w[(2M+3)/2M]^{1/2}$ at every $\Lambda$, with a shift of order $M^{-1/2}$ at finite $\Lambda$: a delay for $\Lambda=\sqrt5$ and an advance for $\Lambda=-\sqrt5$. This is consistent with the ray picture, in which the packet spends the round trip almost entirely at $m\ge2$, where the rows of the ladder matrix do not depend on $\Lambda$. The shift of the nodes near $u=0$ described by (\ref{tr:shift}) varies across the slices that make up the packet, and the spacing at $u=0$ alone does not set the return time.

The comparison of section~\ref{tr:reflsec} excluded 41 of 360 cases, in which the revival was less than 20 times the non-recurrent response at $0.75\,T_{\mathrm{rec}}$: 39 at $\Lambda=\sqrt5$, where that response is the Landau-damped direct response (mostly $M\le64$), and 2 at $\Lambda=\infty$ with $X\ge15$. In the other 319 cases the relative error of the exponent has a median of about $4/M$ and a maximum of about $6.5/M$ for $n=2$ to $6$, at every $X$ tested (up to 30) and every damping per Hermite step $\nu/[|k|s(2M)^{1/2}]$ up to $0.3$. At $\Lambda=\infty$ it comes from the departures of $(2n+1)\mathsf C_{jj}/\nu$ and of $T_{\mathrm{rec}}/2t_w$ from 1, both of order $1/M$. At $\Lambda=\sqrt5$ the revival time and the damping of the slices near $u=0$ depart from their $\Lambda=\infty$ values by more than order $1/M$ (see the paragraph above and section~\ref{tr:collsec}), and we have not determined how these departures combine in the exponent. The law overestimates $\mathcal R$ in all 319 cases. The law applies to a linear phase-mixing problem, or a problem along frozen field lines, with a reflecting closure. It does not apply with an absorbing closure at $m=M$, and the revival is not visible when $e^{-X}$ falls below the non-recurrent response.

Under white-noise forcing of $g_0$, packets injected at different times are uncorrelated, so the steady spectrum at $m$ is the sum of the energies that the injected packets carry through $m$ on all their passes. In the ray picture the damping accumulated between $m=0$ and $m$ is the fraction $\mu=(m/M)^{n+1/2}$ of that accumulated on the way to $M$, because $\int_0^{t}\nu(t'/t_w)^{2n}\,\dd t'$ with $t=t_w(m/M)^{1/2}$ equals $\mu\,\nu t_w/(2n+1)$. The forward pass at $m$ has therefore lost the factor $e^{-X\mu}$ in energy, the backward pass after reflection the factor $e^{-X(2-\mu)}$, and each further round trip multiplies both by $e^{-2X}$. Summing the geometric series, and dividing by the spectrum without reflection, which consists of the first forward pass only, gives (\ref{tr:coth}). For $M=128$, $n=6$, $\nu=3$, $\Lambda=\sqrt5$ and $k_zs=2\pi$, the spectrum without reflection lies 4 to 9\% below the asymptote $\epsilon_W/[s|k_z|(2m)^{1/2}]$ of section~\ref{sec:decorrelation} at $m=2$ and 4, and within $1.6\%$ of it for $16\le m\le64$. With the zero closure the spectrum lies 13 to 19\% above that asymptote for $4\le m\le112$. For $m\le64$ the factor $\coth X=1.185$, together with the deficit of up to 4\% of the open spectrum, accounts for this. At $m=128$ the spectrum lies 27\% above the asymptote; there the closure adds a boundary layer that (\ref{tr:coth}) does not describe. With the ray value of $X$ the law would give an excess of 21\% instead of 18.5\%.
